% answers_short.tex —— 2015 年·福建卷·理科 · 简短答案速查
% 由 tools/make_short_answers.py 自 answers.tex 的 \gksolhead/\gkanswer 提取，
% 并与 2015/2015_数学真题_福建理_答案要点.md 逐题比对后生成，勿手改（改源 MD 后重跑）。

\begin{gknotice}
\textbf{使用说明}：本节只给最终答案便于快速核对，完整推导见解析版。答案为编写者独立推导，\textbf{非官方标准答案}。
\end{gknotice}

\gkpart{选择题}
\begin{center}\begin{tabular}{ccccccccccc}
\toprule
题号 & 1 & 2 & 3 & 4 & 5 & 6 & 7 & 8 & 9 & 10 \\
\midrule
答案 & C & D & B & B & B & C & B & D & A & C \\
\bottomrule
\end{tabular}\end{center}

\gkpart{填空题}
\begin{itemize}[leftmargin=2.6em]
  \item[\textbf{11.}] $80$
  \item[\textbf{12.}] $7$
  \item[\textbf{13.}] $\dfrac5{12}$
  \item[\textbf{14.}] $1<a\leqslant2$
  \item[\textbf{15.}] $5$
\end{itemize}

\gkpart{解答题}
\begin{itemize}[leftmargin=2.6em]
  \item[\textbf{16.}] (1) $\dfrac12$；(2) $P(X=1)=\dfrac16$，$P(X=2)=\dfrac16$，$P(X=3)=\dfrac23$；数学期望 $E(X)=\dfrac52$
  \item[\textbf{17.}] (1) 结论成立，$GF\parallel$ 平面 $ADE$；(2) $\dfrac23$
  \item[\textbf{18.}] (1) $\dfrac{x^2}{4}+\dfrac{y^2}{2}=1$；(2) 点 $G$ 对任意 $m\in\mathbb{R}$ 均在以 $AB$ 为直径的圆外部
  \item[\textbf{19.}] (1) $f(x)=2\sin x$；对称轴为 $x=\dfrac\pi2+k\pi$（$k\in\mathbb{Z}$）；(2) ① $-\sqrt5<m<\sqrt5$；② $\cos(\alpha-\beta)=\dfrac{2m^2}{5}-1$
  \item[\textbf{20.}] (1) 结论成立；(2) 结论成立；(3) $k=1$
  \item[\textbf{21.}] (1) $A^{-1}=\begin{pmatrix}\dfrac32&-\dfrac12\\-2&1\end{pmatrix}$；(2) $C=\begin{pmatrix}\dfrac32&2\\-2&-3\end{pmatrix}$
  \item[\textbf{22.}] (1) 圆 $C$ 的普通方程为 $(x-1)^2+(y+2)^2=9$；直线 $l$ 的直角坐标方程为 $y-x=m$；(2) $m=-3\pm2\sqrt2$
  \item[\textbf{23.}] (1) $a+b+c=4$；(2) 最小值为 $\dfrac87$
\end{itemize}
